\section{政体 dossier：东汉——雒阳复接与地方化}
\label{sec:polity-eastern-han}

东汉把政治中枢放在雒阳，靠近关东和河北的道路。复兴的成就不在于汉的名号重新出现，而在于军队、旧吏、城邑、簿籍和粮道能再度为同一个朝廷服务。度田、裁兵、赈贷和地方任命是重建这种关系的尝试；诏令可说明目标，不能据此推定每一郡县已完成同样的清查或减免。

明、章时期仍须把资源分给北边、西域、羌地和宫廷。窦宪北伐、班超经营绿洲和甘英西使证明雒阳能把人和信息送向远方；它们也依赖河西、使节、驻军与地方王权，不能被写成恒定疆界。和帝诛窦氏后，宦官进入守门、传诏与宫廷安全的关键位置，外戚与宦官随后的轮替正来自这些职能无法由皇帝单独完成。

太学、察举、石经、师友与地方评价，把知识和声望接到官署，也让谁有资格评价他人成为竞争。党锢破坏了这种通道的一部分；压制不能自动让地方精英消失，只会改变他们与州郡、宗族和军队的关系。与此同时，羌地和北方边郡的军费、迁徙和兵员网络持续积累，雒阳留下的账目远少于实际付出的成本。

黄巾之后，镇压和救济需要多处自筹兵粮。州刺史、州牧、边将与豪强并非在某日取代朝廷，而是分别握住征发、任命、仓储、道路或名士网络。董卓、袁绍、曹操、孙氏和刘备等集团正从这些已分开的接口取得资源。献帝仍能提供诏令与名分，却无法使各区域回到同一优先次序；二二〇年的禅代为这种分裂提供仪式形式，而非其唯一原因。

\begin{turningpoint}[东汉政体的得失]
东汉成功重建了可跨区域调度的朝廷，并在边郡、学术和行政上保留了广阔的活动半径；它未能使宫门、边军、财政和地方声望受同一套可持续的检验。危机时为救急而加强的中介，最终成为区域竞争者的资源。
\end{turningpoint}
